\documentclass[8pt]{beamer}

\usetheme{Copenhagen}
\usecolortheme{beaver}
\usepackage[small,center]{caption}
\usepackage{times}
\usefonttheme{structurebold}
\usepackage[english]{babel}
\usepackage{pgf,pgfarrows,pgfnodes,pgfautomata,pgfheaps}
\usepackage{amsmath,amssymb,amsthm}

\usepackage{caption}
\usepackage{tikz}
\usetikzlibrary{shapes.misc}

\newcommand*{\vcenterimage}[1]{\vcenter{\hbox{\includegraphics[width=2in]{#1}}}}
\newcommand*{\vcenterarrow}{\vcenter{\hbox{$\Longrightarrow$}}}

\DeclareMathOperator{\hyphen}{-}

% \definecolor{RPIred}{rgb}{ 0.87,0.12, 0.20}
\definecolor{ballblue}{rgb}{0.13, 0.67, 0.8}
\definecolor{lightgray}{rgb}{0.83, 0.83, 0.83}
%\setbeamercolor{block title}{bg=lightgray,fg=RPIred}
\setbeamercolor{block body}{bg=white,fg=black}
\setbeamercovered{dynamic}
%\setbeamercolor*{item}{fg=RPIred}

\DeclareMathOperator*{\argmin}{arg\,min}

\captionsetup[subfigure]{labelformat=empty}
\captionsetup[figure]{labelformat=empty}
\setbeamertemplate{navigation symbols}{}
\setbeamertemplate{footline}[frame number]
\begin{document}

\include{notation}




\def\signed #1{{\leavevmode\unskip\nobreak\hfil\penalty50\hskip2em
  \hbox{}\nobreak\hfil #1 %
  \parfillskip=0pt \finalhyphendemerits=0 \endgraf}}

\newsavebox\mybox
\newenvironment{aquote}[1]
  {\savebox\mybox{#1}\begin{quote}}
  {\signed{\usebox\mybox}\end{quote}}



% tikz stuff
\tikzset{cross/.style={cross out, draw=black, minimum size=2*(#1-\pgflinewidth),
inner sep=0pt, outer sep=0pt},
%default radius will be 1pt.
cross/.default={2.5pt}}



\frame{
\title{\Large Cardinal: A framework for doing FSI with deal.II + IBAMR}

\author{{\Large David Wells \\\vspace{0.1in} University of North Carolina, Chapel Hill}\\
\vspace{0.2in} {In collaboration with:\\{}Laryssa Abdala, Marshall Davey, Simone Rossi, Boyce Griffith}}

\date{August 12 2024\\{} deal.II Workshop}

\begin{figure}[h]
\centering
%\includegraphics[width=1.5in]{RPI_letterhead.pdf}
 \end{figure}%

\vspace{-0.2in}
\titlepage
}

\section{Introduction to Cardinal}
\begin{frame}
    \frametitle{Governing Equations}
    Uses the immersed-boundary finite element / finite difference method (IFED) (Boffi and Heltai)
    \begin{itemize}
      \item \emph{Eulerian} finite difference method for Navier-Stokes
      \item \emph{Lagrangian} finite element method for mechanics
    \end{itemize}

    \begin{align*}
        \rho \frac{D\ub}{Dt}(\xb, t) &= -\nabla p(\xb, t)
        + \mu \nabla ^2 \ub(\xb, t) + \fb(\xb, t),
        \\
        \nabla \cdot \ub(\xb, t) &= \, 0
        \\
        % no PPs yet
        \fb(\xb, t) &= \int_{\soliddomO} \Fb(\Xb,t) \, \delta(\xb - \Chib (\Xb,t)) \dXb
        \\
        \frac{\partial \Chib}{\partial t}(\Xb, t) &= \ \Ub(\Xb,t) = \int_{\Omega}
        \ub(\xb, t) \, \delta(\xb - \Chib (\Xb,t)) \dxb
        \\
        \int_{\soliddomO} \Fb(\Xb,t) \cdot \psib(\Xb) \dXb &=
        -\int_{\soliddomO} \PPs(\Xb,t) : \nabla_{\Xb} \psib(\Xb) \dXb
    \end{align*}

    \begin{center}
      The structure's velocity $\frac{\partial \Chib}{\partial t}$ is \emph{interpolated} from the fluid $\ub$; the force on the fluid $\fb$ is \emph{spread} from the structure $\Fb$.
    \end{center}
\end{frame}

\begin{frame}
    \frametitle{Regularized Delta Functions} Heart of the IB method: we replace
    $\delta(\xb)$ with $\delta_h(\xb)$. Here, $\delta_h(\xb)$ has a width of $2 \Delta x$.

    \begin{figure}
      \centering
      \begin{tabular}{c c}
        \includegraphics[width=0.4\linewidth]{interactionstencils/elemental_stencil.pdf} &
        \includegraphics[width=0.4\linewidth]{interactionstencils/nodal_stencil.pdf}
      \end{tabular}

      Can either evaluate the integral operators with Gauss quadrature (left) or
      nodal quadrature (right). Not shown to scale!
    \end{figure}

    IBAMR implements FDM and IB methods - use deal.II for all unstructured grid things.
\end{frame}

\begin{frame}
    \frametitle{Goals}
    %%% TODO: picture!
    What is this project all about? Haven't people already done everything with the IB method?
    \begin{itemize}
      \item Immersed methods (not cut cells), neutral buoyancy. IFED now, IIM soon.
      \item Modern fiber-reinforced material models.
      \item All features are available through input files.
      \item Parallel (FEBio has more features but will probably never support MPI).
      \item Students can do interesting things with it and \emph{not} write 5000 line monoliths
    \end{itemize}
\end{frame}

\begin{frame}
    \frametitle{Brief History}
    How did we get here?
    \begin{itemize}
      \item 2018: four-chambered-heart. Realistic cardiac dynamics with libMesh!
      \item 2020: deal.II with simplices is ``not impossible'' (what did we write on the wiki?)
      \item 2021: four-chambered-heart, IBAMR, deal.II (fiddle)
      \item 2022: I write another libMesh-based paper
      \item 2023: Heart model (Davey) plus IFED (me) plus electrocardiology (Abdala, Rossi)
      \item 2024: Realistic valve dynamics with deal.II
    \end{itemize}
\end{frame}

\section{Rest of the talk}
\begin{frame}
  \frametitle{What is this talk about?}
  \begin{itemize}
    \item All the problems Laryssa and I ran into
    \item Interesting Cardinal features
    \item Ideas for improving other deal.II apps
  \end{itemize}

  \vfill

  \begin{center}
    Framing: eliminating the three kinds of waste (muda, mura, muri).
  \end{center}
\end{frame}

\section{Who are we?}
\begin{frame}
    \frametitle{Who are we?}
    We are...
    \begin{itemize}
      \item Faculty members, research scientists, postdocs
      \item Decades of C++ experience
      \item Formal training in mathematics, computer science
    \end{itemize}

    Our jobs are...
    \begin{itemize}
      \item dev ops (who is going to fix the CI?)
      \item write proposals
      \item teach classes
    \end{itemize}
\end{frame}

\begin{frame}
    \frametitle{Who are graduate students?}
    \begin{aquote}{W. Edwards Deming}
      Put a good person in a bad system and the bad system wins, no contest.
    \end{aquote}

    Students are...
    \begin{itemize}
      \item Smart
      \item Know a lot about subject areas (biomechanics, electrophysiology)
      \item New to scientific computing
    \end{itemize}

    \vfill
    Students do not have...
    \begin{itemize}
      \item Years of C++ experience
      \item Years of engineering experience
      \item Formal training in computer science
    \end{itemize}
\end{frame}

\section{Muda}
\begin{frame}
  \frametitle{Muda - non-value-added work}
  \begin{aquote}{Sigeo Shingo}
    The most dangerous kind of waste is the waste we do not recognize.
  \end{aquote}
  What do we actually do that is of value?

  \pause

  \begin{itemize}
    \item Graduate students
    \item Write papers
    \item Ship software
  \end{itemize}

  Everything else is either necessary (e.g., write grants) or unnecessary
  (should be removed).
\end{frame}

\begin{frame}
  \frametitle{non-value-added work}

  Some non-value-added work. If we watch our students - what are they doing?
  \begin{itemize}
    \item Compiling code (\emph{fixes}: ccache, CI, better hardware, use \texttt{-O1})
    \item Debugging (\emph{don't write bugs!})
    \item Running everything in debug mode (too slow!)
    \item Managing dependencies, upgrades, system administration
    \item Typing too slowly or inaccurately
  \end{itemize}
\end{frame}

\begin{frame}
  \frametitle{non-value-added work: compiling}

  Various performance improvements are not obvious - students don't have the vocabulary.
  \begin{itemize}
    \item \texttt{make} to \texttt{ninja} (what is \texttt{-j4}?)
    \item \texttt{-O3} to \texttt{-O1} (multiple builds?)
    \item \texttt{ccache} (what are object files?)
    \item Make the CI fail fast, finish in 5 minutes
    \item Trickery to make \texttt{-DDEBUG} work with \texttt{-O1} and not \texttt{-g}
  \end{itemize}
\end{frame}

\section{Mura}

\section{Muri}

\end{document}
